\chapter{DX-004 --- Smooth Surface and String Composition Seams}\label{record-336}

\noindent{\footnotesize\color{BookMuted}\path{docs/method/retro/DX_004-smooth-surface-and-string-composition-seams.md}\par}\smallskip

\begin{RecordMeta}
title: DX-004 — Smooth Surface and String Composition Seams
lane: retro
legend: DX

\end{RecordMeta}

\section{Disposition}\label{record-336-disposition}

Shipped a public mixed composition seam in @flyingrobots/bijou-tui: contentSurface() now bridges raw text into surface-land, vstackSurface()/hstackSurface()/placeSurface() accept string \textbar{} Surface inputs at the composition edge, runtime view normalization points callers at the sanctioned bridge, and package docs now explain how to keep composition on the surface side without loosening App.view back to raw strings.

\section{Original Proposal}\label{record-336-original-proposal}

Legend: \href{https://github.com/flyingrobots/bijou/blob/1853ad902f6bf219326409e7f743a5a1d6517dbc/docs/method/legends/DX-developer-experience.md}{DX --- Developer Experience}

\section{Idea}\label{record-336-idea}

Reduce confusion at the boundary between string-oriented composition and \texttt{S\allowbreak{}u\allowbreak{}r\allowbreak{}f\allowbreak{}a\allowbreak{}c\allowbreak{}e}/layout-oriented composition.

\section{Why}\label{record-336-why}

Bijou currently exposes a real conceptual seam here, but the ergonomics are easy to trip over:

\begin{itemize}
\tightlist
\item
  some helpers return \texttt{S\allowbreak{}u\allowbreak{}r\allowbreak{}f\allowbreak{}a\allowbreak{}c\allowbreak{}e}
\item
  some stack/layout helpers still expect strings
\item
  \texttt{v\allowbreak{}i\allowbreak{}e\allowbreak{}w\allowbreak{}(\allowbreak{})} requires \texttt{S\allowbreak{}u\allowbreak{}r\allowbreak{}f\allowbreak{}a\allowbreak{}c\allowbreak{}e\allowbreak{} \allowbreak{}|\allowbreak{} \allowbreak{}L\allowbreak{}a\allowbreak{}y\allowbreak{}o\allowbreak{}u\allowbreak{}t\allowbreak{}N\allowbreak{}o\allowbreak{}d\allowbreak{}e}, not string
\item
  migration from older string-first patterns is easy to misunderstand
\end{itemize}

The issue is not only documentation. Some composition helpers may need better bridging or clearer contracts.

\section{Captured feedback}\label{record-336-captured-feedback}

From \texttt{w\allowbreak{}a\allowbreak{}r\allowbreak{}p\allowbreak{}-\allowbreak{}t\allowbreak{}t\allowbreak{}d} TUI development:

\begin{itemize}
\tightlist
\item
  \texttt{b\allowbreak{}a\allowbreak{}d\allowbreak{}g\allowbreak{}e\allowbreak{}(\allowbreak{})} returning \texttt{S\allowbreak{}u\allowbreak{}r\allowbreak{}f\allowbreak{}a\allowbreak{}c\allowbreak{}e} made \texttt{v\allowbreak{}s\allowbreak{}t\allowbreak{}a\allowbreak{}c\allowbreak{}k\allowbreak{}(\allowbreak{})} composition awkward
\item
  \texttt{s\allowbreak{}u\allowbreak{}r\allowbreak{}f\allowbreak{}a\allowbreak{}c\allowbreak{}e\allowbreak{}T\allowbreak{}o\allowbreak{}S\allowbreak{}t\allowbreak{}r\allowbreak{}i\allowbreak{}n\allowbreak{}g\allowbreak{}(\allowbreak{}.\allowbreak{}.\allowbreak{}.\allowbreak{})} was required in some places but forbidden in \texttt{v\allowbreak{}i\allowbreak{}e\allowbreak{}w\allowbreak{}(\allowbreak{})}
\item
  the \texttt{v\allowbreak{}i\allowbreak{}e\allowbreak{}w\allowbreak{}(\allowbreak{})} return constraint is understandable but still easy to trip over during migration
\end{itemize}

\section{Likely scope}\label{record-336-likely-scope}

\begin{itemize}
\tightlist
\item
  audit string-land vs surface-land composition boundaries
\item
  decide where bridging helpers should exist and where they should not
\item
  improve the composition story for small surface-returning primitives
\item
  tighten docs or error guidance where the seam is intentionally strict
\end{itemize}

